\MajorSection{ملاحظات على الترجمة}

\MinorSection{الناموس الأعلى}

تُرجم \textenglish{Higher Law} بـ«الناموس الأعلى»، وتُرجم \textenglish{Law} بـ«الناموس» حين يدل على النظام الذي يسير به الكون أو تتحدد به الأفعال. أما في المواضع الدينية التي يدل فيها على أحكام الكتاب أو حرف الشريعة، فاستُعملت «الشريعة». ويشرح برتون نفسه، في الملاحظة الأولى، أن الناموس لا يفترض عند الصوفي واضعًا شخصيًا له. و«أنشودة» مقابل \textenglish{Lay} تسمية أدبية، لا تحديد لوزن عربي.

\MinorSection{الزاهد والصوفي}

يشرح برتون الزاهد بوصفه المتمسك بحرف الشريعة، والصوفي بوصفه المتمسك بروحها. فالزاهد هنا شخصية في الجدل، لا مجرد إنسان قليل المتاع. وكذلك تتسع صورة الصوفي لديه لتمثل العارف الشاك وصابّ الخمر والمجوسي العجوز. وقد بقيت هذه الأدوار المختلفة ظاهرة بدل دمجها في لفظ واحد.

\MinorSection{العقل والذهن والدماغ}

استُعمل «العقل» أساسًا لـ\textenglish{Reason}، و«الذهن» حيث يلزم إبراز \textenglish{Mind} في مقابل الدماغ، و«الدماغ» لـ\textenglish{Brain}، و«الفكر» لـ\textenglish{Thought}. وينطلق شرح برتون من اتصال القوى العقلية بفعل الدماغ والعصب، ثم يميز بين ما يعجز العقل عن معرفته وما يناقض العقل. وهذه من الفروق التي تحكم القسم السابع.

\MinorSection{الروح والنفس}

نُقل \textenglish{Soul} في معظم الجدل عن الخلود بـ«الروح»، إذ يناقش برتون فكرة كيان يبقى بعد الجسد. وحيث يجمع \textenglish{Spirit} و\textenglish{Soul} و\textenglish{Shade}، حُفظ تعدد الألفاظ بـ«الروح والنفس والظل». أما \textenglish{Self} و\textenglish{I} فيردان، بحسب الموضع، «الذات» و«النفس» و«أنا». وفي الملاحظة الأولى يصرّح برتون بأنه يفهم الروح حالًا من الأحوال وكلمة للإحساس بالهوية الفردية؛ ثم يتناول العواطف بوصفها حقائق فاعلة في الحياة.

\MinorSection{الرجولة وتنمية النفس}

حُفظت «الرجولة» لـ\textenglish{Manhood} في النداء الأخلاقي، كما حُفظت دلالتها على النضج في سياق العمر. ويجمع النص بين الشجاعة والاستقلال وطلب الكمال الفردي، مع المعجم الذكوري لزمانه. و«تنمية النفس» مقابل \textenglish{Self cultivation} تشمل في هذا السياق التعلم وتكوين الذات، مع مراعاة الآخرين؛ ويشرح برتون امتزاج «أنا» بـ«نحن» في الملاحظة الأولى.

\MinorSection{الأجناس والأسماء التاريخية}

تختلف دلالة \textenglish{Race} بحسب السياق بين الجنس البشري والعرق والشعب والسلالة. وتظهر الآريّة والطورانيّة والقبائل «الزنوجية» ضمن تصنيفات برتون نفسه. أما «الموري» في القسم الخامس فقد حُفظ اسمًا تاريخيًا؛ لأن برتون يشرح أنه يقصد به الدانكالي وغيرهم من القبائل التي يسمّيها بهذا الوصف، فلا يصح اختزاله هنا إلى «المغربي». و«الفرنجي» في حوار القسم الثامن مسيحي أوروبي كما تحدده ملاحظته. و«المجوس» تقابل \textenglish{Guebres} أو \textenglish{Magians} بحسب الموضع؛ و«الملتشا» اسم يشرح برتون أنه يُطلق على غير الهندوس.

\MinorSection{ذيل الذئب وإشارات المطلع}

«الملكة» في المطلع صورة للقمر المتناقص. و«ذيل الذئب» صورة الفجر الكاذب، الذي يشرحه برتون بضوء الفجر البروجي. وحُفظت صورة الذيل في القصيدة، ثم أُثبت تفسيرها في حاشيته. و«الأزل والأبد» في القسم الثاني ينقلان وجهي ما يسميه النص \textenglish{two Eternities}، كما يشرحهما في الملاحظة الثانية.

\MinorSection{المادة والحياة في لغة القرن التاسع عشر}

حُفظت «البروتوبلازم» و«بقعة الهلام» و«سلسلة الوجود»، مع صورة السلّم وتدرّج القوى والكائنات في القصيدة. وترد أسماء أسلاف الحصان والتصنيفات الجيولوجية والتفسير الجمجمي للقوى العقلية كما يوردها برتون. وفي الملاحظة الأولى نُقل \textenglish{Radiant Matter} بـ«المادة المتألقة»؛ فالعبارة تنتمي إلى مناقشات الحالة الرابعة للمادة وتجارب التفريغ الكهربائي في أنابيب مفرّغة، كما في محاضرة وليام كروكس سنة ١٨٧٩. والمقارنة تفسر المصطلح؛ أما ربطه بإمكان عالم الأرواح في الفقرة نفسها فهو استنتاج الكاتب الذي يقتبسه برتون.

\MinorSection{النيرفانا والتصورات الدينية}

يعرّف برتون النيرفانا بـ«الفناء النسبي»، ثم بـ«الانطفاء الجزئي بالاندماج في الأعلى»، ويميّزها من البارينيرفانا التي يسمّيها «الفناء المطلق». نُقلت تعريفاته هذه بألفاظها، وكذلك تفسيره للكارما والتناسخ. وتعود إليه أيضًا مساواة الخضر بإيليا، وروايته عن رجم الحلاج، ومقارناته بين موسى والصوفي والظاهري والباطني. تُقرأ هذه الأقوال ضمن تصوره للأديان وتاريخه المقارن.

\MinorSection{الجدل في الحقيقة والإيمان}

يحفظ القسم السادس الفرق بين الحقيقة الموضوعية والحقيقة الأخلاقية، وتشرح الملاحظة الثانية أن القضية الكاذبة تظل قضية موجودة مثل القضية الصادقة؛ فوجود القضية غير صدقها. وحُفظ أيضًا التوتر بين «كل إيمان باطل، وكل إيمان حق» وبين الدعوة إلى تعليق الحكم. ونُقل \textenglish{un know} في ختام القسم بـ«يتخلّص مما يعرف» إبرازًا لفعل فكّ المعرفة المألوفة؛ ويمكن مراجعته لاحقًا إلى «يعرف كيف لا يعرف» بحسب الإيقاع المختار.

\MinorSection{الأرقام والاقتباسات}

أُبقيت الأرقام والنسب والتواريخ ومراجع الآيات والأبيات كما وردت في النص المعتمد؛ ومنها نسبتا المسيحيين والمسلمين وتأريخ الأسر المصرية. وتُرجمت الاقتباسات اللاتينية والفرنسية والإيطالية والإسبانية والإنجليزية القديمة إلى العربية، مع إثبات بعض الألفاظ الأصلية اللازمة لفهم الجدل. وفي وصف الحور، حُفظ شرح برتون ولفظه «أحور»؛ وصيغة المؤنث العربية «حوراء». والمراجعة التحقيقية للاقتباسات ومصادرها مرحلة مستقلة عن هذه الترجمة الأولية.
